% Curs 9, pagina 7 din 8 — transcriere a manuscrisului din 1994.
% Fragment destinat adăugării după pagina 6 în courses/course-09.tex.
% Linia „Obs. Orb x = ... = Gx” tăiată de autor nu este reprodusă;
% sunt păstrate observaţiile netăiate imediat următoare.
% Niciun proiect original nu a fost modificat.

\section*{Curs 9 - Pagina 7}

\noindent\underline{Prop.} Fie $G$ gr. şi $A$ o $G$-mulţ. Rel. binară pe $A$
\[
x\sim_G y\ \overset{\mathrm{def}}{\Longleftrightarrow}
(\exists)\,a\in G\text{ a.î. }a x=y,\qquad x,y\in A.
\]
\noindent (rel. de echiv.)

\bigskip
\noindent Clasele de echiv. ale ac. rel. s.n. \underline{orbitele acţ. lui $G$ pe $A$}.

\bigskip
\noindent\underline{Dem.} Arăt că:
\begin{align*}
&x\sim_G x,\qquad (\forall)\,x\in A,\\
&x\sim_G y\ \Longrightarrow\ y\sim_G x,\\
&x\sim_G y\text{ şi }y\sim_G z\ \Longrightarrow\ x\sim_G z.
\end{align*}
\noindent Reflexivitatea rezultă din (2) (de la def. $G$-mulţimii).

\bigskip
\noindent $(\exists)\,a\in G$ a.î. $a x=y$ $\Longrightarrow$
\[
x=e x=(a^{-1}a)x=a^{-1}(a x)=a^{-1}y.
\]

\bigskip
\noindent $(\exists)\,a,b\in G$ a.î.
\[
\left.\begin{aligned}
 a x&=y,\\
 b y&=z
\end{aligned}\right\}
\quad\Longrightarrow\quad z=b(a x)=(ba)x.
\]

\bigskip
\noindent\underline{Obs.} Dacă $x\in A$, $\operatorname{Orb}(x)=\operatorname{Orb}_G(x)$
(orbita lui $x$ prin acţ. lui $G$).
\[
\operatorname{Orb}(x)=\{a x\mid a\in G\}=Gx.
\]
\[
\operatorname{Orb}(x)=\widehat{x}
=\{y\in A\mid y\sim_G x\}.
\]

\bigskip
\noindent\underline{Def.} Dacă $A$ e o $G$-mulţ. şi $x\in A$,
\[
\operatorname{Stab}_G(x)\ \overset{\mathrm{def}}{=}
\{a\in G\mid a\cdot x=x\}
\]
\noindent stabilizatorul lui $x$ (în $G$).

\noindent Evident $e\in\operatorname{Stab}_G(x)$.

\bigskip
\noindent\fbox{T} Dacă $A$ e o $G$-mulţ. şi $x\in A$, atunci:
\begin{enumerate}
\item[(1)] $\operatorname{Stab}_G(x)\leq G$.
\item[(2)] $\operatorname{card}(\operatorname{Orb}(x))
=[G:\operatorname{Stab}_G(x)]$.
\item[(3)] $G$ finit $\Longrightarrow\operatorname{card}(\operatorname{Orb}(x))\mid\operatorname{ord}G$.
\end{enumerate}

\bigskip
\noindent $S_n\times A\longrightarrow A$, \qquad $A=\{1,2,\ldots,n\}$,
\[
\operatorname{Stab}_{S_n}(n)
=\{\sigma\in S_n\mid\sigma(n)=n\}.
\]
\noindent $\sigma\cdot n=\sigma(n)\in A$.
